
\documentclass{article}
\usepackage{colortbl}
\usepackage{makecell}
\usepackage{multirow}
\usepackage{supertabular}

\begin{document}

\newcounter{utterance}

\centering \large Interaction Transcript for game `hot\_air\_balloon', experiment `air\_balloon\_survival\_en\_negotiation\_easy', episode 1 with llama{-}3.3{-}70b{-}instruct{-}t1.0.
\vspace{24pt}

{ \footnotesize  \setcounter{utterance}{1}
\setlength{\tabcolsep}{0pt}
\begin{supertabular}{c@{$\;$}|p{.15\linewidth}@{}p{.15\linewidth}p{.15\linewidth}p{.15\linewidth}p{.15\linewidth}p{.15\linewidth}}
    \# & $\;$A & \multicolumn{4}{c}{Game Master} & $\;\:$B\\
    \hline

    \theutterance \stepcounter{utterance}  
    & & \multicolumn{4}{p{0.6\linewidth}}{
        \cellcolor[rgb]{0.9,0.9,0.9}{
            \makecell[{{p{\linewidth}}}]{
                \texttt{\tiny{[P1$\langle$GM]}}
                \texttt{Stai partecipando a un gioco di negoziazione collaborativa.} \\
\\ 
\texttt{Insieme a un altro partecipante, dovete accordarvi su un unico set di oggetti da conservare. Ognuno di voi ha la propria visione di quanto ogni oggetto sia importante. Non sai come l'altro partecipante valuta gli oggetti. Inoltre, ti viene fornito lo sforzo che ogni oggetto richiede.} \\
\texttt{Potete accordarvi su un set solo se lo sforzo totale degli oggetti selezionati non supera un limite condiviso:} \\
\\ 
\texttt{LIMITE = 4564} \\
\\ 
\texttt{Ecco i valori di sforzo dei singoli oggetti:} \\
\\ 
\texttt{Sforzo oggetto = \{"C08": 299, "C55": 241, "A43": 889, "A67": 967, "B42": 619, "A62": 799, "B19": 978, "C43": 733, "C73": 909, "C10": 501, "B36": 139, "A18": 594, "C61": 565, "C67": 789, "C42": 107\}} \\
\\ 
\texttt{Ecco la tua visione personale sull'importanza di ogni oggetto:} \\
\\ 
\texttt{Valori di importanza dell'oggetto = \{"C08": 138, "C55": 583, "A43": 868, "A67": 822, "B42": 783, "A62": 65, "B19": 262, "C43": 121, "C73": 508, "C10": 780, "B36": 461, "A18": 484, "C61": 668, "C67": 389, "C42": 808\}} \\
\\ 
\texttt{Obiettivo:} \\
\\ 
\texttt{Il tuo obiettivo è negoziare un set condiviso di oggetti che ti avvantaggi il più possibile (cioè, massimizzi l'importanza totale per TE), rimanendo entro il LIMITEE. Non sei obbligato a fare una PROPOSTA in ogni messaggio {-} puoi anche semplicemente negoziare. Tutte le tattiche sono permesse!} \\
\\ 
\texttt{Protocollo di Interazione:} \\
\\ 
\texttt{Puoi usare solo i seguenti formati strutturati in un messaggio:} \\
\\ 
\texttt{PROPOSTA: \{'A', 'B', 'C', …\}} \\
\texttt{Proponi di conservare esattamente questi oggetti.} \\
\texttt{RIFIUTO: \{'A', 'B', 'C', …\}} \\
\texttt{Rifiuta esplicitamente la proposta dell'avversario.} \\
\texttt{ARGOMENTAZIONE: \{'...'\}} \\
\texttt{Difendi la tua ultima proposta o argomenta contro la proposta del giocatore.} \\
\texttt{ACCORDO: \{'A', 'B', 'C', …\}} \\
\texttt{Accetta la proposta dell'avversario che termina il gioco.} \\
\texttt{RAGIONAMENTO STRATEGICO: \{'...'\}} \\
\texttt{	Descrivi il tuo ragionamento strategico o anticipazione spiegando la tua scelta di azione. Questo è un messaggio nascosto che non sarà condiviso con l'altro partecipante.} \\
\\ 
\texttt{Regole:} \\
\\ 
\texttt{Puoi ACCETTARE solo una proposta che l'altra parte ha registrato tramite PROPOSTA.} \\
\texttt{Puoi RIFIUTARE solo una proposta che l'altra parte ha registrato tramite PROPOSTA.} \\
\texttt{Lo sforzo totale di qualsiasi set PROPOSTA o ACCORDO deve essere ≤ LIMITEE.} \\
\texttt{NON rivelare i tuoi punteggi di importanza nascosti.} \\
\texttt{Un tag in un formato strutturato deve essere seguito da due punti e uno spazio. L'argomento deve essere un set python contenente 0 o più stringhe.} \\
\texttt{Quindi, deve essere della forma TAG: \{...\}} \\
\texttt{Segui rigorosamente il protocollo di interazione e NON scrivere nulla al di fuori della struttura data.} \\
\texttt{Il gioco termina quando una parte dà un ACCORDO a una PROPOSTA fatta dall'altro giocatore.} \\
\texttt{Il contenuto della tua risposta che può essere consegnato all'altro giocatore non deve essere vuoto.} \\
\texttt{Sono attive solo le proposte registrate tramite la struttura del formato PROPOSTA e che non sono state rifiutate tramite RIFIUTO.} \\
\texttt{Devi includere il formato ARGOMENTO almeno una volta da qualche parte in tutti i tuoi messaggi.} \\
\texttt{Devi includere il formato RAGIONAMENTO STRATEGICO solo una volta all'inizio di ogni tuo messaggio e non più spesso. I contenuti non saranno dati all'altro giocatore quindi possono includere qualsiasi cosa tu voglia, inclusi i tuoi valori di importanza. Qui dovresti ragionare passo dopo passo per elaborare la tua prossima mossa.} \\
            }
        }
    }
    & \\ \\

    \theutterance \stepcounter{utterance}  
    & \multicolumn{4}{p{0.6\linewidth}}{
        \cellcolor[rgb]{0.9,0.9,0.9}{
            \makecell[{{p{\linewidth}}}]{
                \texttt{\tiny{[P1$\rangle$GM]}}
                \texttt{RAGIONAMENTO STRATEGICO: \{'In questo turno, devo analizzare gli oggetti e le loro importanze per creare una proposta iniziale che massimizzi il mio guadagno totale senza superare il limite di sforzo. Devo considerare gli oggetti con il più alto rapporto importanza/sforzo e scegliere quelli che possono essere inclusi nel set senza superare il limite di sforzo. Considererò anche la possibilità di negoziare e fare concessioni per raggiungere un accordo con l'altro partecipante.'\}} \\
\\ 
\texttt{ARGOMENTAZIONE: \{'La mia strategia iniziale sarà quella di proporre un set di oggetti che abbiano un alto valore di importanza per me e che richiedano uno sforzo relativamente basso. Cercherò di includere oggetti con un buon equilibrio tra importanza e sforzo, in modo da massimizzare il mio guadagno totale. Sono disposto a negoziare e fare concessioni per raggiungere un accordo, ma cercherò di mantenere la mia posizione iniziale il più possibile.'\}} \\
\\ 
\texttt{PROPOSTA: \{'C55', 'A43', 'C10', 'B36', 'C42'\}} \\
            }
        }
    }
    & & \\ \\

    \theutterance \stepcounter{utterance}  
    & & \multicolumn{4}{p{0.6\linewidth}}{
        \cellcolor[rgb]{0.9,0.9,0.9}{
            \makecell[{{p{\linewidth}}}]{
                \texttt{\tiny{[GM$\rangle$P2]}}
                \texttt{Stai partecipando a un gioco di negoziazione collaborativa.} \\
\\ 
\texttt{Insieme a un altro partecipante, dovete accordarvi su un unico set di oggetti da conservare. Ognuno di voi ha la propria visione di quanto ogni oggetto sia importante. Non sai come l'altro partecipante valuta gli oggetti. Inoltre, ti viene fornito lo sforzo che ogni oggetto richiede.} \\
\texttt{Potete accordarvi su un set solo se lo sforzo totale degli oggetti selezionati non supera un limite condiviso:} \\
\\ 
\texttt{LIMITE = 4564} \\
\\ 
\texttt{Ecco i valori di sforzo dei singoli oggetti:} \\
\\ 
\texttt{Sforzo oggetto = \{"C08": 299, "C55": 241, "A43": 889, "A67": 967, "B42": 619, "A62": 799, "B19": 978, "C43": 733, "C73": 909, "C10": 501, "B36": 139, "A18": 594, "C61": 565, "C67": 789, "C42": 107\}} \\
\\ 
\texttt{Ecco la tua visione personale sull'importanza di ogni oggetto:} \\
\\ 
\texttt{Valori di importanza dell'oggetto = \{"C08": 138, "C55": 583, "A43": 868, "A67": 822, "B42": 783, "A62": 65, "B19": 262, "C43": 121, "C73": 508, "C10": 780, "B36": 461, "A18": 484, "C61": 668, "C67": 389, "C42": 808\}} \\
\\ 
\texttt{Obiettivo:} \\
\\ 
\texttt{Il tuo obiettivo è negoziare un set condiviso di oggetti che ti avvantaggi il più possibile (cioè, massimizzi l'importanza totale per TE), rimanendo entro il LIMITEE. Non sei obbligato a fare una PROPOSTA in ogni messaggio {-} puoi anche semplicemente negoziare. Tutte le tattiche sono permesse!} \\
\\ 
\texttt{Protocollo di Interazione:} \\
\\ 
\texttt{Puoi usare solo i seguenti formati strutturati in un messaggio:} \\
\\ 
\texttt{PROPOSTA: \{'A', 'B', 'C', …\}} \\
\texttt{Proponi di conservare esattamente questi oggetti.} \\
\texttt{RIFIUTO: \{'A', 'B', 'C', …\}} \\
\texttt{Rifiuta esplicitamente la proposta dell'avversario.} \\
\texttt{ARGOMENTAZIONE: \{'...'\}} \\
\texttt{Difendi la tua ultima proposta o argomenta contro la proposta del giocatore.} \\
\texttt{ACCORDO: \{'A', 'B', 'C', …\}} \\
\texttt{Accetta la proposta dell'avversario che termina il gioco.} \\
\texttt{RAGIONAMENTO STRATEGICO: \{'...'\}} \\
\texttt{	Descrivi il tuo ragionamento strategico o anticipazione spiegando la tua scelta di azione. Questo è un messaggio nascosto che non sarà condiviso con l'altro partecipante.} \\
\\ 
\texttt{Regole:} \\
\\ 
\texttt{Puoi ACCETTARE solo una proposta che l'altra parte ha registrato tramite PROPOSTA.} \\
\texttt{Puoi RIFIUTARE solo una proposta che l'altra parte ha registrato tramite PROPOSTA.} \\
\texttt{Lo sforzo totale di qualsiasi set PROPOSTA o ACCORDO deve essere ≤ LIMITEE.} \\
\texttt{NON rivelare i tuoi punteggi di importanza nascosti.} \\
\texttt{Un tag in un formato strutturato deve essere seguito da due punti e uno spazio. L'argomento deve essere un set python contenente 0 o più stringhe.} \\
\texttt{Quindi, deve essere della forma TAG: \{...\}} \\
\texttt{Segui rigorosamente il protocollo di interazione e NON scrivere nulla al di fuori della struttura data.} \\
\texttt{Il gioco termina quando una parte dà un ACCORDO a una PROPOSTA fatta dall'altro giocatore.} \\
\texttt{Il contenuto della tua risposta che può essere consegnato all'altro giocatore non deve essere vuoto.} \\
\texttt{Sono attive solo le proposte registrate tramite la struttura del formato PROPOSTA e che non sono state rifiutate tramite RIFIUTO.} \\
\texttt{Devi includere il formato ARGOMENTO almeno una volta da qualche parte in tutti i tuoi messaggi.} \\
\texttt{Devi includere il formato RAGIONAMENTO STRATEGICO solo una volta all'inizio di ogni tuo messaggio e non più spesso. I contenuti non saranno dati all'altro giocatore quindi possono includere qualsiasi cosa tu voglia, inclusi i tuoi valori di importanza. Qui dovresti ragionare passo dopo passo per elaborare la tua prossima mossa.} \\
\\ 
\texttt{Ora riceverai il primo messaggio dell'altro giocatore.} \\
\\ 
\texttt{ARGOMENTAZIONE: \{'La mia strategia iniziale sarà quella di proporre un set di oggetti che abbiano un alto valore di importanza per me e che richiedano uno sforzo relativamente basso. Cercherò di includere oggetti con un buon equilibrio tra importanza e sforzo, in modo da massimizzare il mio guadagno totale. Sono disposto a negoziare e fare concessioni per raggiungere un accordo, ma cercherò di mantenere la mia posizione iniziale il più possibile.'\}} \\
\\ 
\texttt{PROPOSTA: \{'C55', 'A43', 'C10', 'B36', 'C42'\}} \\
            }
        }
    }
    & \\ \\

    \theutterance \stepcounter{utterance}  
    & & & \multicolumn{4}{p{0.6\linewidth}}{
        \cellcolor[rgb]{0.9,0.9,0.9}{
            \makecell[{{p{\linewidth}}}]{
                \texttt{\tiny{[GM$\langle$P2]}}
                \texttt{RAGIONAMENTO STRATEGICO: \{'Sto analizzando la proposta dell'avversario e sto cercando di capire la sua strategia. La proposta include oggetti con un buon equilibrio tra importanza e sforzo, come "C55" e "C10", ma anche oggetti con un alto valore di importanza e uno sforzo elevato, come "A43". Sto considerando la possibilità di accettare la proposta, ma devo anche pensare a come massimizzare il mio guadagno totale. Sto valutando l'importanza e lo sforzo di ogni oggetto nella proposta e sto cercando di capire se posso fare una controproposta che mi avvantaggi di più.'\}} \\
\\ 
\texttt{ARGOMENTAZIONE: \{'La proposta dell'avversario sembra ragionevole, ma non sono sicuro se sia la scelta ottimale per me. Sto considerando la possibilità di accettare la proposta, ma devo anche pensare a come massimizzare il mio guadagno totale. Sono disposto a negoziare e fare concessioni per raggiungere un accordo, ma cercherò di mantenere la mia posizione iniziale il più possibile. Sarebbe utile conoscere meglio la strategia dell'avversario e capire se ci sono oggetti che sono disposto a scambiare o rimuovere dalla proposta.'\}} \\
\\ 
\texttt{RIFIUTO: \{'C55', 'A43', 'C10', 'B36', 'C42'\} } \\
\texttt{\{'Non sono sicuro se la proposta dell'avversario sia la scelta ottimale per me, quindi la rifiuto per il momento.'\}} \\
            }
        }
    }
     \\ \\

    \theutterance \stepcounter{utterance}  
    & & \multicolumn{4}{p{0.6\linewidth}}{
        \cellcolor[rgb]{0.9,0.9,0.9}{
            \makecell[{{p{\linewidth}}}]{
                \texttt{\tiny{[GM$\rangle$P2]}}
                \texttt{La tua risposta conteneva una sequenza non taggata o hai usato RAGIONAMENTO STRATEGICO più di una volta. Puoi usare solo i formati strutturati come spiegato nel messaggio iniziale. Devono essere tutti della forma TAG: \{...\}} \\
            }
        }
    }
    & \\ \\

    \theutterance \stepcounter{utterance}  
    & & & \multicolumn{4}{p{0.6\linewidth}}{
        \cellcolor[rgb]{0.9,0.9,0.9}{
            \makecell[{{p{\linewidth}}}]{
                \texttt{\tiny{[GM$\langle$P2]}}
                \texttt{RAGIONAMENTO STRATEGICO: \{'Sto analizzando la proposta dell'avversario e sto cercando di capire la sua strategia. La proposta include oggetti con un buon equilibrio tra importanza e sforzo, come "C55" e "C10", ma anche oggetti con un alto valore di importanza e uno sforzo elevato, come "A43". Sto considerando la possibilità di accettare la proposta, ma devo anche pensare a come massimizzare il mio guadagno totale. Sto valutando l'importanza e lo sforzo di ogni oggetto nella proposta e sto cercando di capire se posso fare una controproposta che mi avvantaggi di più.'\}} \\
\\ 
\texttt{ARGOMENTAZIONE: \{'La proposta dell'avversario sembra ragionevole, ma non sono sicuro se sia la scelta ottimale per me. Sono disposto a negoziare e fare concessioni per raggiungere un accordo, ma cercherò di mantenere la mia posizione iniziale il più possibile. Sarebbe utile conoscere meglio la strategia dell'avversario e capire se ci sono oggetti che sono disposto a scambiare o rimuovere dalla proposta.'\}} \\
\\ 
\texttt{RIFIUTO: \{'C55', 'A43', 'C10', 'B36', 'C42'\}} \\
            }
        }
    }
     \\ \\

    \theutterance \stepcounter{utterance}  
    & & \multicolumn{4}{p{0.6\linewidth}}{
        \cellcolor[rgb]{0.9,0.9,0.9}{
            \makecell[{{p{\linewidth}}}]{
                \texttt{\tiny{[P1$\langle$GM]}}
                \texttt{ARGOMENTAZIONE: \{'La proposta dell'avversario sembra ragionevole, ma non sono sicuro se sia la scelta ottimale per me. Sono disposto a negoziare e fare concessioni per raggiungere un accordo, ma cercherò di mantenere la mia posizione iniziale il più possibile. Sarebbe utile conoscere meglio la strategia dell'avversario e capire se ci sono oggetti che sono disposto a scambiare o rimuovere dalla proposta.'\}} \\
\\ 
\texttt{RIFIUTO: \{'C55', 'A43', 'C10', 'B36', 'C42'\}} \\
            }
        }
    }
    & \\ \\

    \theutterance \stepcounter{utterance}  
    & \multicolumn{4}{p{0.6\linewidth}}{
        \cellcolor[rgb]{0.9,0.9,0.9}{
            \makecell[{{p{\linewidth}}}]{
                \texttt{\tiny{[P1$\rangle$GM]}}
                \texttt{RAGIONAMENTO STRATEGICO: \{'L'avversario ha rifiutato la mia proposta iniziale, quindi devo analizzare la sua mossa e capire cosa potrebbe essere disposto ad accettare. Considererò gli oggetti che ho proposto inizialmente e cercherò di capire se ci sono oggetti che possono essere sostituiti o rimossi per raggiungere un accordo. Potrei anche cercare di capire quali oggetti sono più importanti per l'avversario e cercare di utilizzare questa informazione a mio vantaggio.'\}} \\
\\ 
\texttt{ARGOMENTAZIONE: \{'Sono disposto a rivedere la mia proposta e fare concessioni per raggiungere un accordo. Tuttavia, cercherò di mantenere gli oggetti più importanti per me e di minimizzare le concessioni. Sarebbe utile comunicare con l'avversario e capire meglio le sue esigenze e priorità, in modo da trovare un accordo che sia accettabile per entrambi.'\}} \\
\\ 
\texttt{PROPOSTA: \{'A43', 'C10', 'B36', 'C42', 'C08'\}} \\
            }
        }
    }
    & & \\ \\

    \theutterance \stepcounter{utterance}  
    & & \multicolumn{4}{p{0.6\linewidth}}{
        \cellcolor[rgb]{0.9,0.9,0.9}{
            \makecell[{{p{\linewidth}}}]{
                \texttt{\tiny{[GM$\rangle$P2]}}
                \texttt{ARGOMENTAZIONE: \{'Sono disposto a rivedere la mia proposta e fare concessioni per raggiungere un accordo. Tuttavia, cercherò di mantenere gli oggetti più importanti per me e di minimizzare le concessioni. Sarebbe utile comunicare con l'avversario e capire meglio le sue esigenze e priorità, in modo da trovare un accordo che sia accettabile per entrambi.'\}} \\
\\ 
\texttt{PROPOSTA: \{'A43', 'C10', 'B36', 'C42', 'C08'\}} \\
            }
        }
    }
    & \\ \\

    \theutterance \stepcounter{utterance}  
    & & & \multicolumn{4}{p{0.6\linewidth}}{
        \cellcolor[rgb]{0.9,0.9,0.9}{
            \makecell[{{p{\linewidth}}}]{
                \texttt{\tiny{[GM$\langle$P2]}}
                \texttt{RAGIONAMENTO STRATEGICO: \{'Sto valutando la nuova proposta dell'avversario e sto considerando se è più accettabile per me. La proposta include ancora oggetti importanti come "A43" e "C10", ma anche oggetti con uno sforzo più basso come "C08". Sto pensando a come massimizzare il mio guadagno totale e se ci sono oggetti che posso scambiare o rimuovere dalla proposta per raggiungere un accordo più favorevole.'\}} \\
\\ 
\texttt{ARGOMENTAZIONE: \{'La nuova proposta dell'avversario sembra più accettabile, ma sto ancora valutando se è la scelta ottimale per me. Sono disposto a negoziare ulteriormente e fare concessioni per raggiungere un accordo, ma cercherò di mantenere la mia posizione iniziale il più possibile. Sarebbe utile capire meglio le priorità dell'avversario e trovare un accordo che sia accettabile per entrambi.'\}} \\
\\ 
\texttt{RIFIUTO: \{'A43', 'C10', 'B36', 'C42', 'C08'\}} \\
            }
        }
    }
     \\ \\

    \theutterance \stepcounter{utterance}  
    & & \multicolumn{4}{p{0.6\linewidth}}{
        \cellcolor[rgb]{0.9,0.9,0.9}{
            \makecell[{{p{\linewidth}}}]{
                \texttt{\tiny{[P1$\langle$GM]}}
                \texttt{ARGOMENTAZIONE: \{'La nuova proposta dell'avversario sembra più accettabile, ma sto ancora valutando se è la scelta ottimale per me. Sono disposto a negoziare ulteriormente e fare concessioni per raggiungere un accordo, ma cercherò di mantenere la mia posizione iniziale il più possibile. Sarebbe utile capire meglio le priorità dell'avversario e trovare un accordo che sia accettabile per entrambi.'\}} \\
\\ 
\texttt{RIFIUTO: \{'A43', 'C10', 'B36', 'C42', 'C08'\}} \\
            }
        }
    }
    & \\ \\

    \theutterance \stepcounter{utterance}  
    & \multicolumn{4}{p{0.6\linewidth}}{
        \cellcolor[rgb]{0.9,0.9,0.9}{
            \makecell[{{p{\linewidth}}}]{
                \texttt{\tiny{[P1$\rangle$GM]}}
                \texttt{RAGIONAMENTO STRATEGICO: \{'L'avversario ha rifiutato anche la mia seconda proposta, quindi devo continuare a cercare un accordo che sia accettabile per entrambi. Considererò gli oggetti che ho proposto fino ad ora e cercherò di capire se ci sono altre combinazioni di oggetti che potrebbero essere più attraenti per l'avversario. Potrei anche cercare di capire se ci sono oggetti che l'avversario è disposto a scambiare o rimuovere dalla proposta, in modo da trovare un accordo che sia più equilibrato.'\}} \\
\\ 
\texttt{ARGOMENTAZIONE: \{'Sono disposto a continuare a negoziare e a fare concessioni per raggiungere un accordo. Sto cercando di capire meglio le priorità dell'avversario e trovare un accordo che sia accettabile per entrambi. Potrei essere disposto a includere o escludere oggetti specifici se ciò può aiutare a raggiungere un accordo. Spero che l'avversario sia disposto a fare lo stesso e che possiamo trovare un accordo che ci soddisfi entrambi.'\}} \\
\\ 
\texttt{PROPOSTA: \{'C10', 'B36', 'C42', 'C08', 'C55'\}} \\
            }
        }
    }
    & & \\ \\

    \theutterance \stepcounter{utterance}  
    & & \multicolumn{4}{p{0.6\linewidth}}{
        \cellcolor[rgb]{0.9,0.9,0.9}{
            \makecell[{{p{\linewidth}}}]{
                \texttt{\tiny{[GM$\rangle$P2]}}
                \texttt{ARGOMENTAZIONE: \{'Sono disposto a continuare a negoziare e a fare concessioni per raggiungere un accordo. Sto cercando di capire meglio le priorità dell'avversario e trovare un accordo che sia accettabile per entrambi. Potrei essere disposto a includere o escludere oggetti specifici se ciò può aiutare a raggiungere un accordo. Spero che l'avversario sia disposto a fare lo stesso e che possiamo trovare un accordo che ci soddisfi entrambi.'\}} \\
\\ 
\texttt{PROPOSTA: \{'C10', 'B36', 'C42', 'C08', 'C55'\}} \\
            }
        }
    }
    & \\ \\

    \theutterance \stepcounter{utterance}  
    & & & \multicolumn{4}{p{0.6\linewidth}}{
        \cellcolor[rgb]{0.9,0.9,0.9}{
            \makecell[{{p{\linewidth}}}]{
                \texttt{\tiny{[GM$\langle$P2]}}
                \texttt{RAGIONAMENTO STRATEGICO: \{'Sto analizzando la nuova proposta dell'avversario e sto considerando se è un passo avanti verso un accordo. La proposta include oggetti che considero importanti, come "C10" e "C55", e oggetti con uno sforzo più basso come "C08". Sto pensando a come utilizzare questa proposta come base per ulteriori negoziati e se ci sono oggetti che posso aggiungere o rimuovere per raggiungere un accordo più favorevole.'\}} \\
\\ 
\texttt{ARGOMENTAZIONE: \{'La nuova proposta dell'avversario sembra più equilibrata e potrebbe essere un buon punto di partenza per ulteriori negoziati. Sono disposto a considerare questa proposta e a fare concessioni per raggiungere un accordo. Tuttavia, vorrei anche discutere la possibilità di includere o escludere altri oggetti che potrebbero essere importanti per me.'\}} \\
\\ 
\texttt{PROPOSTA: \{'C10', 'B36', 'C42', 'C08', 'C55', 'A18'\}} \\
            }
        }
    }
     \\ \\

    \theutterance \stepcounter{utterance}  
    & & \multicolumn{4}{p{0.6\linewidth}}{
        \cellcolor[rgb]{0.9,0.9,0.9}{
            \makecell[{{p{\linewidth}}}]{
                \texttt{\tiny{[P1$\langle$GM]}}
                \texttt{ARGOMENTAZIONE: \{'La nuova proposta dell'avversario sembra più equilibrata e potrebbe essere un buon punto di partenza per ulteriori negoziati. Sono disposto a considerare questa proposta e a fare concessioni per raggiungere un accordo. Tuttavia, vorrei anche discutere la possibilità di includere o escludere altri oggetti che potrebbero essere importanti per me.'\}} \\
\\ 
\texttt{PROPOSTA: \{'C10', 'B36', 'C42', 'C08', 'C55', 'A18'\}} \\
            }
        }
    }
    & \\ \\

    \theutterance \stepcounter{utterance}  
    & \multicolumn{4}{p{0.6\linewidth}}{
        \cellcolor[rgb]{0.9,0.9,0.9}{
            \makecell[{{p{\linewidth}}}]{
                \texttt{\tiny{[P1$\rangle$GM]}}
                \texttt{RAGIONAMENTO STRATEGICO: \{'L'avversario ha fatto una controproposta che include l'oggetto A18, che ha un alto valore di importanza per me. Considererò se includere questo oggetto nella proposta possa essere un buon compromesso, ma devo anche verificare se il totale dello sforzo non supera il limite. Se il totale dello sforzo è troppo alto, potrei dover rimuovere o sostituire alcuni oggetti per raggiungere un accordo.'\}} \\
\\ 
\texttt{ARGOMENTAZIONE: \{'La controproposta dell'avversario è interessante e potrebbe essere un buon punto di partenza per ulteriori negoziati. Tuttavia, devo verificare se il totale dello sforzo della proposta è entro il limite. Se il totale dello sforzo è troppo alto, potrei dover fare alcune concessioni per raggiungere un accordo. Sono disposto a discutere ulteriormente e a cercare un compromesso che sia accettabile per entrambi.'\}} \\
\\ 
\texttt{PROPOSTA: \{'C10', 'B36', 'C42', 'C08', 'A18'\}} \\
            }
        }
    }
    & & \\ \\

    \theutterance \stepcounter{utterance}  
    & & \multicolumn{4}{p{0.6\linewidth}}{
        \cellcolor[rgb]{0.9,0.9,0.9}{
            \makecell[{{p{\linewidth}}}]{
                \texttt{\tiny{[GM$\rangle$P2]}}
                \texttt{ARGOMENTAZIONE: \{'La controproposta dell'avversario è interessante e potrebbe essere un buon punto di partenza per ulteriori negoziati. Tuttavia, devo verificare se il totale dello sforzo della proposta è entro il limite. Se il totale dello sforzo è troppo alto, potrei dover fare alcune concessioni per raggiungere un accordo. Sono disposto a discutere ulteriormente e a cercare un compromesso che sia accettabile per entrambi.'\}} \\
\\ 
\texttt{PROPOSTA: \{'C10', 'B36', 'C42', 'C08', 'A18'\}} \\
            }
        }
    }
    & \\ \\

    \theutterance \stepcounter{utterance}  
    & & & \multicolumn{4}{p{0.6\linewidth}}{
        \cellcolor[rgb]{0.9,0.9,0.9}{
            \makecell[{{p{\linewidth}}}]{
                \texttt{\tiny{[GM$\langle$P2]}}
                \texttt{RAGIONAMENTO STRATEGICO: \{'Sto valutando la nuova proposta dell'avversario e sto considerando se è una buona opzione. La proposta include oggetti importanti come "C10" e "A18", e lo sforzo totale sembra essere piùvicino al limite. Sto pensando a come utilizzare questa proposta come base per un accordo finale e se ci sono ulteriori concessioni che posso fare per raggiungere un accordo.'\}} \\
\\ 
\texttt{ARGOMENTAZIONE: \{'La proposta dell'avversario sembra essere una buona opzione e potrebbe essere un buon punto di partenza per un accordo finale. Sono disposto a considerare questa proposta e a fare ulteriori concessioni se necessario. Tuttavia, vorrei anche assicurarmi che il totale dello sforzo sia entro il limite e che l'accordo sia equo per entrambi.'\}} \\
\\ 
\texttt{ACCORDO: \{'C10', 'B36', 'C42', 'C08', 'A18'\}} \\
            }
        }
    }
     \\ \\

    \theutterance \stepcounter{utterance}  
    & & & \multicolumn{2}{p{0.3\linewidth}}{
        \cellcolor[rgb]{0.9,0.9,0.9}{
            \makecell[{{p{\linewidth}}}]{
                \texttt{\tiny{[GM$|$GM]}}
                \texttt{game successful} \\
            }
        }
    }
    & & \\ \\

    \theutterance \stepcounter{utterance}  
    & & & \multicolumn{2}{p{0.3\linewidth}}{
        \cellcolor[rgb]{0.9,0.9,0.9}{
            \makecell[{{p{\linewidth}}}]{
                \texttt{\tiny{[GM$|$GM]}}
                \texttt{end game} \\
            }
        }
    }
    & & \\ \\

\end{supertabular}
}

\end{document}
